% ===================================================================== PROTOCOL
\section{The \sage{} protocol}
\label{sec:protocol}

\sage{} combines a single-finalizer dual-run, a quorum-certified boundary, an $n-f$-signed
migration manifest, evidence of durable source retirement, and deadline-bounded rollback under
two-tier finality. \Cref{alg:sage} specifies the design-level controller for a correct validator.

\subsection{Migration schedule as a transition system}
Governance fixes a schedule $\hd<\hc\le\hr$ in platform state $P$; the schedule can be amended only
through the legacy engine before $\hd$. Here $\hc$ is the \emph{earliest} admissible transition; the
\CutCert{} later fixes the actual height $\hs=\hb+1\ge\hc$, and a reversible handoff requires
$\hs<\hr$. Authority is defined by the transition system of \cref{fig:lts} over the phases
$\{\phase{v1},\phase{dual},\phase{v2},\phase{rollback},\phase{sealed}\}$:
\begin{equation}
  \Auth(h)=
  \begin{cases}
    \Csrc^{g_1} & \text{in }\phase{v1}\text{ and }\phase{dual},\\
    \Csrc^{g_1'},\ g_1'>g_1 & \text{in }\phase{rollback},\\
    \Ctgt & \text{in }\phase{v2}\text{ and }\phase{sealed}.
  \end{cases}
  \label{eq:authority}
\end{equation}
Here $\Auth$ names the designated engine generation, not the existence of a usable finalization
quorum. After a valid \CutCert{} forms, durable Rule~\ref{rule:lock} and $q_1>2f$ exclude source
finalization above $\hb$, while the target is not yet enabled. This certification pause has no
permitted finalizer above the boundary, even though the controller remains in \phase{dual}.
Abort has priority over every other transition in \phase{v2}, and \phase{dual} self-loops until the
readiness predicate holds. The component state machine splits \phase{dual} into \phase{fencing} and
\phase{activating} sub-phases, which order manifest adoption, durable fence installation, and
certificate collection ahead of target bootstrap in the specified contract. The complete network
path through these stages is not established by the component states alone.

\begin{figure*}[t]
\centering
\includegraphics[width=\textwidth]{figures/fig_lts.pdf}
\caption{\sage{} as a labelled transition system. Target authority requires the \CutCert, an adopted
\ManCert, a durable local fence, and a \FenceCert. In \phase{v2} a valid \AbortCert{} before $\hr$
takes priority and binds a fresh generation $g_1'>g_1$. After $\hr$, only a \SealCert{} reaches
\phase{sealed}.}
\label{fig:lts}
\end{figure*}

\begin{figure*}[t]
\centering
\includegraphics[width=\textwidth]{figures/fig_arch.pdf}
\caption{Dual-run data path. The legacy engine alone finalizes while the target re-executes finalized
blocks; target authority requires \CutCert, \ManCert, and \FenceCert{} stages.}
\label{fig:arch}
\end{figure*}

\subsection{Single-finalizer dual-run}
The central invariant is that at every height at most one engine is authorized to finalize. During
$[\hd,\hs)$, $\Csrc$ orders and finalizes $B_h$ exactly as before $\hd$, and $\Ctgt$ never proposes.

\begin{definition}[Shadow mode]
\label{def:shadow}
An engine runs in shadow mode at height $h$ if, for the canonical block $B_h$ finalized by the
authoritative engine, each correct validator runs the shadow engine's validation on $B_h$. It records
$\sigma_h\in\{\mathsf{valid},\mathsf{invalid}\}$ and a root $\tilde r_h$ (with $\tilde r_h=\bot$ on
rejection), and does not propose or finalize. Shadow validation is a side-effect-free function of
$(B_h,X_{h-1})$ and never blocks finalization.
\end{definition}

\Cref{fig:arch} shows the resulting data path. The controller of validator $i$ collects the shadow
verdicts, and the validator is \emph{ready} at height $h$ when the last $\kappa$ verdicts are valid and
the roots match:
\begin{equation}
  \begin{aligned}
  \mathrm{Ready}_i(h)\iff{}& h\ge\hc-1\\
  &\wedge\;\forall j\in(h-\kappa,\,h]:\ \sigma_j=\mathsf{valid}\wedge\tilde r_j=r_j .
  \end{aligned}
  \label{eq:ready}
\end{equation}

\subsection{Quorum-certified cutover}
Governance cannot precommit a hash that does not yet exist. It therefore fixes only the policy part
$G=\langle\mathrm{cid},e,\hc,g_1,g_2\rangle$ before dual-run. The concrete boundary tuple depends on
the data:
\begin{equation}
  \beta=\bigl\langle \mathrm{cid},\,e,\,\hb,\,\Hash(B_{\hb}),\,r_{\hb}\bigr\rangle,\qquad \hs=\hb+1 .
  \label{eq:beta}
\end{equation}

\begin{definition}[Readiness attestation and \CutCert]
\label{def:cutcert}
In epoch $e$, a correct validator durably locks its first eligible $\beta$ and emits at most one
attestation $\mathsf{att}_i=\mathrm{Sign}_i\langle\textsc{ready},\beta\rangle$. It signs only when
$\mathrm{Ready}_i(\hb)$ holds and it has locally finalized and recomputed $B_{\hb}$. A \CutCert{}
$\pi_{\mathrm B}$ is a set of $n-f$ valid attestations on one identical tuple. Attestations that name
different heights, parents, or roots never combine. If the committee does not converge on one body,
the migration stays pending. Selecting a later candidate would require a separate close-and-retry
protocol, which this paper does not analyze.
\end{definition}

\begin{rulebox}[Attestation lock]
\label{rule:lock}
A correct validator may attest at $\hb$ only if it has issued no source proposal, vote, or
finalization authorization above $\hb$ in generation $g_1$, including authorizations prepared for
transmission. Before releasing $\mathsf{att}_i$, it atomically persists both its one-body slot
$(e,\beta)$ and the prohibition of every such authorization at $h>\hb$ in $g_1$. Attestation and
source-signing paths serialize against this durable record: neither buffered work nor an alternate
source handler may bypass it. The write completes before the signature becomes externally visible.
A restart reloads and enforces the prohibition before enabling signing, retransmission, or consensus
message processing. Ambiguous recovery fails closed. The prohibition never expires in $g_1$;
certified recovery may enable only a fresh $g_1'>g_1$.
\end{rulebox}

Rule~\ref{rule:lock} makes a certified boundary actionable. Without its prohibition, source
authorizations above $\hb$ issued before local fence installation may disqualify correct signers
from later fence sharing and leave the handoff pending (\cref{lem:fence-live}). The
fence-only theorem does not assume Rule~\ref{rule:lock}; the admitted locked design does, and its
\CutCert{} already excludes a source quorum under $q_1>2f$ (\cref{lem:fence-live}). Persisting
$\beta$ without enforcing the prohibition after restart is insufficient. This complete invariant
has not been established on the evaluated network path (\cref{tab:status}).

The cost is a deliberate source pause during certification. \Cref{th:liveness} bounds it only on
the successful, post-GST path with a common eligible candidate. If validators lock incompatible
candidates, or too few eligible signers remain, they cannot unlock by timeout or silently switch to
a later height. Migration may then remain pending indefinitely; a safe close-and-retry protocol is
outside this design. This is a practical liveness restriction, not just a missing optimization.

The implementation carries $\pi_{\mathrm B}$ as an $n-f$ Ed25519 multi-signature, that is, a list of
individual EUF-CMA signatures with size $O(n)$. A threshold or Boneh--Lynn--Shacham (BLS) aggregate
could preserve the abstract threshold argument, but would require a compatible signer/key model,
setup and share validation. No aggregated implementation or timing is evaluated here.

\subsection{Quorum-signed migration manifest}
The complete signing body is
\begin{equation}
  \mathcal M_{\mathrm B}=\bigl\langle \mathrm{cid},e,\mathrm{cfg},\hb{+}1,\Hash(B_{\hb}),r_{\hb},
  g_1,g_2,\Hash_{\mathrm{tool}},r_{\mathrm{ctx}},\Hash(\pi_{\mathrm B})\bigr\rangle ,
  \label{eq:manifest}
\end{equation}
where $\Hash_{\mathrm{tool}}$ binds the migration tooling and $r_{\mathrm{ctx}}$ is the optional
replay-context root. A \ManCert{} $\pi_{\mathrm M}$ contains $n-f$ shares over $\Hash(\mathcal M_{\mathrm
B})$ in the manifest authority domain. A correct validator accepts only if (i) $\pi_{\mathrm B}$
verifies on $\beta$; (ii) $\pi_{\mathrm M}$ verifies on the complete body; (iii) the boundary parent and
root match its own finalized legacy tip; and (iv) its durable epoch slot is empty or already holds
the same body. Each correct validator issues at most one manifest share per epoch, and the durable
slot enforces this.

\subsection{Source-fence certificate}
\begin{definition}[Source fence and \FenceCert]
\label{def:fencecert}
Governance admits a source-finality policy only if its quorum $q_1$ satisfies
\eqref{eq:intro-fence}. For a certified boundary with transition height $\hs$, a correct validator
(1)~durably adopts the complete manifest; (2)~durably installs a fence that rejects source proposals,
votes, and finalization at every $h\ge\hs$; (3)~confirms that it has issued no source authorization
at any $h\ge\hs$; and only then (4)~releases a domain-separated fence share. The share binds the
chain, epoch, source engine and generation, $\hs$, boundary block and root, $\Hash(\pi_{\mathrm B})$,
$\Hash(\mathcal M_{\mathrm B})$, the source-finality policy, and the fence policy. A \FenceCert{}
contains $q_F$ matching shares. Target authority is exposed only after the validator verifies the
\ManCert{} and the \FenceCert{} and confirms its own durable fence. A correct validator never releases
a fence share before installation and never resumes a fenced generation.
\end{definition}

\begin{figure*}[t]
\centering
\includegraphics[width=\textwidth]{figures/fig_seq.pdf}
\caption{Certified handoff at a correct validator. Manifest adoption and durable fence installation
precede fence sharing; the \FenceCert{} precedes target authority. Dashed dual-run messages never finalize.}
\label{fig:seq}
\end{figure*}

\Cref{fig:seq} shows the specified order as messages at one validator. The component contract
requires verification and persistence of the \CutCert, durable manifest adoption, and a fence write
before sharing. Once $q_F$ matching shares are present, the contract requires assembly and
verification of the \FenceCert{} against the committee and the admitted quorum pair, followed by
persistence. Target authority may be exposed only when the durable fence, adopted manifest, and
verified \FenceCert{} agree on the same boundary. A conforming restart reloads those records before
target bootstrap. Both governance admission and \FenceCert{} verification evaluate
\eqref{eq:intro-fence}. These component operations do not constitute an exercised network
transport path for the three certificates (\cref{tab:status}). The extra fence round certifies
policy-bound retirement and manifest adoption by its correct signers; it does not further reduce
the residual voting set already excluded by a valid \CutCert{} under durable \cref{rule:lock}.

\paragraph{State and proof dependencies}
The two branches below distinguish retirement from activation:
\begin{equation}
\begin{gathered}
\text{durable Rule~\ref{rule:lock}}+\CutCert+(q_1>2f)\\
\Longrightarrow\ \text{source quiescence},\\[0.3ex]
\CutCert\to\ManCert\to\text{local durable fence}\\
\to\text{fence share}\to\FenceCert,\\[0.3ex]
\CutCert+\ManCert+\text{local durable fence}+\FenceCert\\
\Longrightarrow\ \text{activation permitted}.
\end{gathered}
\label{eq:dependencies}
\end{equation}
The first branch proves old-generation exclusion under the durable-lock premise. The second
certifies manifest- and policy-bound retirement by its correct signers; the last is the chosen
per-validator activation gate, not a proof that the certificate is independently necessary for
source exclusion. The fence-only theorem uses the second branch without the readiness-lock
premise. No branch assumes that every validator receives the certificates.

\subsection{Two-tier finality and bounded rollback}
Legacy blocks below $\hs$ are \emph{absolutely} final. Target blocks are \emph{provisional} in
\phase{v2}; they become absolutely final only when a valid $n-f$ \SealCert{} moves the controller to
\phase{sealed}. The deadline $\hr$ closes the rollback window and makes sealing eligible, but height
alone never seals. A client receipt reports the tier that the committing engine stamped into the block
header, and that tier sits inside the signed finality payload. A serving validator therefore cannot
relabel a provisional block as absolute without invalidating the certificate. When a \SealCert{}
covers a prefix, its proof can establish effective absolute finality for receipts at or below the
sealed height; the original signed provisional receipt must remain unchanged. Height alone is not
a promotion proof. End-to-end client receipt reconciliation is outside the evaluated path.
Cross-chain actions and other irreversible external effects must wait for a verified \SealCert{}
covering the relevant block. Restoring ledger state cannot revoke an external payment, message, or
physical action that was already performed.

Local and certified aborts have different effects. A target stall or a manifest failure may emit an
AbortVote and restore the local state to the anchor $X_{\hs-1}$, but committee authority does not
change. Only a valid $n-f$ \AbortCert{}, issued before $\hr$ under an exclusive terminal-signing lock
and binding a fresh $g_1'>g_1$, returns authority to $\Csrc^{g_1'}$. The old generation stays fenced.
Timeouts are detected by a progress ticker that runs off the finalization path. It is armed when
target authority is exposed, refreshed only by a strictly advancing target commit, and refused at or
after $\hr$.

\paragraph{Replay context}
Deterministic replay of a discarded suffix onto $\Csrc^{g_1'}$ requires that execution depends only on
the ordered transactions and parent state (\cref{as:metadata}). The manifest binds a replay-context
root. Each provisional block carries the chain identifier, height, timestamp, beneficiary, base fee,
randomness, and an optional oracle snapshot. A missing, reordered, or mismatched record makes rollback
\emph{fail closed}, which turns metadata dependence into an explicit, checkable input. The canonical
record occupies $114+|\mathrm{cid}|$ bytes, or $122$\,B with an eight-byte chain identifier, plus $40$\,B
for an oracle root. A $1000$-block window therefore needs about $119$--$158$\,KiB per validator.
\Cref{tab:rollback-compat} summarizes the behaviour by workload class.

\begin{table}[t]
\centering
\caption{Design-level rollback contract by workload class. Context admission is implemented;
application execution and client reconciliation are not evaluated.}
\label{tab:rollback-compat}
\small
\renewcommand{\arraystretch}{1.14}
\rowcolors{2}{rowtint}{white}
\begin{tabular}{@{}L{0.30\columnwidth}L{0.24\columnwidth}L{0.36\columnwidth}@{}}
\toprule
\textbf{Workload} & \textbf{Replay eligibility} & \textbf{Required behaviour} \\
\midrule
Metadata-independent & By replay lemma & Execute ordered suffix from anchor \\
Reads block metadata & If captured in context & Consume bound context; else refuse \\
Live oracle input & If snapshot root bound & Consume bound snapshot; else refuse \\
Cross-chain actions on provisional state & Not \sage's to replay & Integrator waits for \SealCert \\
\bottomrule
\end{tabular}
\end{table}

\begin{algorithm}[t]
\caption{Design-level \sage{} controller at a correct validator $i$}
\label{alg:sage}
\small
\begin{algorithmic}[1]
\Require schedule $(\hd,\hc,\hr)$, window $\kappa$, admitted $(q_1,q_F)$ with $q_F+q_1>n+f$
\State $\mathit{phase}\gets\phase{v1}$;\ $\hs\gets\bot$;\ $\mathit{streak}\gets0$
\For{each block, certificate, timer, or recovery event}
  \State on recovery reload Rule~\ref{rule:lock} and generation fences before source use or buffered releases
  \State update retained evidence; let $h$ be the local committed tip
  \If{$\mathit{phase}=\phase{v1}\wedge h\ge\hd$} $\mathit{phase}\gets\phase{dual}$ \EndIf
  \If{$\mathit{phase}=\phase{v2}\wedge h<\hr\wedge$ valid \AbortCert{} binds $g_1'>g_1$}
    \State restore $X_{\hs-1}$; $\Auth\gets\Csrc^{g_1'}$; $\mathit{phase}\gets\phase{rollback}$; \textbf{continue}
  \EndIf
  \If{$\mathit{phase}=\phase{dual}$ and event is a newly finalized block $B_h$}
    \State $(\sigma_h,\tilde r_h)\gets\Ctgt.\textsc{ShadowValidate}(B_h)$
    \State $\mathit{streak}\gets(\sigma_h=\mathsf{valid}\wedge\tilde r_h=r_h)\,?\,\mathit{streak}+1:0$
    \If{$\mathrm{Ready}_i(h)$ and slot $(e,\cdot)$ empty and Rule~\ref{rule:lock} eligibility holds}
      \State atomically persist $(e,\beta)$ and source prohibition before releasing $\mathsf{att}_i$
    \EndIf
  \EndIf
  \If{$\mathit{phase}=\phase{dual}$}
    \State when $n-f$ matching attestations exist: verify and persist $\pi_{\mathrm B}(\beta)$; set $\hs\gets\hb+1$
    \State when $\pi_{\mathrm B}$ and local boundary checks pass: issue one durable manifest share
    \State when valid $\pi_{\mathrm M}$ exists: verify and durably adopt $\mathcal M_{\mathrm B}$
    \State after adoption and no prior conflicting source authorization: install durable fence, then release share
    \If{matching valid \CutCert, \ManCert, \FenceCert{} and local durable fence, with $\hs<\hr$}
      \State initialize target from certified anchor; persist activation; $\Auth\gets\Ctgt$; $\mathit{phase}\gets\phase{v2}$
    \EndIf
  \EndIf
  \State timer expiry may request abort; it never changes authority without a valid \AbortCert
  \If{$\mathit{phase}=\phase{v2}\wedge h\ge\hr\wedge$ valid \SealCert} $\mathit{phase}\gets\phase{sealed}$ \EndIf
\EndFor
\end{algorithmic}
\end{algorithm}

Certificate processing and retransmission are independent of new finalizations: once the source
pauses, a block-only callback cannot complete the remaining handoff. Each stage is idempotent on its
durable identity. The pseudocode specifies the event dependencies; it is not a claim that the
network implementation already implements this dispatcher or terminal-certificate convergence.

\subsection{Design rationale}
The mechanisms have overlapping safety effects and distinct evidence roles. A scheduled blind
cutover supplies no committee-certified boundary. The process controls in \cref{sec:rq2} isolate
unsafe target thresholds from unauthorized crossings; they do not establish that every ungated
switch forks. Dual-run with only local authorization lets
partition sides choose different boundaries, which \cref{th:unique} excludes through certificate
intersection and signer locking. Equality with a proposer-supplied manifest hash authenticates
nothing, so \cref{th:unforge} needs a quorum-signed body. Under durable \cref{rule:lock} and
$q_1>2f$, removing only \FenceCert{} does not permit residual source finalization. The rejected
policy in \cref{cor:cross-split} illustrates the fence-only contract's bound, not necessity inside
the admitted locked design. \FenceCert{} instead attests the complete retirement identity after
manifest adoption (\cref{prop:retirement-evidence}). Without a rollback window, a
target failure after cutover becomes irreversible, which is the Fabric~v3 behaviour. Finality gadgets
such as Casper~FFG~\cite{casperffg} and GRANDPA~\cite{grandpa} add a checkpoint predicate over an
\emph{unchanged} block producer. They neither initialize a new engine's metadata
(\cref{def:bootstrap}) nor certify a reversible boundary between quorum families, so they complement
\sage{} rather than replace it.
